% 张三慧第4版第二卷24章原书思考题24.1--24.11
% 主教材扫描PDF物理页312--313。与章末习题24.1--24.16分开计数。
\section*{第24章\quad 光的偏振：思考题}
\paragraph{思考题 24.1}
自然光常被看成两个相互垂直振动分量之和。
为什么这两个分量之间没有固定的相位关系，
而椭圆偏振光却必须有确定的相位差？
\paragraph{思考题 24.2}
一束光可能为完全线偏振、部分线偏振或自然光，
如何使用转动的理想线偏振分析器分辨这三种情形？
\paragraph{思考题 24.3}
一个偏振片的透振轴没有标注方向，
如何用另一个已知偏振的光源或偏振片确定其透振轴？
\paragraph{思考题 24.4}
一束光入射到两透明介质的平面分界面，
只有透射光而观察不到反射光。
这对入射偏振态和入射角有什么可能的要求？
\paragraph{思考题 24.5}
两个线偏振片正交，原本没有透射光。
中间插入第三片后可出现透光，说明第三片取向如何；
哪些取向仍然完全不透光？
\paragraph{思考题 24.6}
Barkla的双散射实验（图24.37）：
X射线先沿水平方向被一个碳块散射到第二碳块，
而第二碳块在某些垂直观察方向没有散射强度。
从电偶极子辐射的角分布分析为什么这支持X射线的电磁横波性质。
\paragraph{思考题 24.7}
对单轴晶体，若入射光在外部沿晶体光轴的方向射到
一个与光轴倾斜的表面，进入晶体后是否一定不发生双折射？
请根据折射后的波法线方向判断。
\paragraph{思考题 24.8}
若一束光可能是线偏振、椭圆偏振或自然光，
怎样联合使用线偏振分析器与四分之一波片鉴别？
\paragraph{思考题 24.9}
一块四分之一波片与两块线偏振片混在一起且没有标记，
如何通过转动、消光或圆偏振检验辨认出三者？
\paragraph{思考题 24.10}
在原图24.27的晶体偏振干涉装置中，
若去掉前偏振片 \(P_1\) 或后分析器 \(P_2\)，
能否保留可观测的干涉条纹？解释正交偏振模式与相干叠加。
\paragraph{思考题 24.11}
在原图24.28的偏振干涉实验里，
若将偏振片 \(P_1\) 移到晶体 \(C\)
和 \(P_2\) 之间，原干涉强度公式是否还成立？
若两偏振片 \(P_1,P_2\) 平行，结果如何？
